\documentclass[11pt,a4paper]{article}

\newcommand{\doctitle}{Threat Model y Cadena de Ataque}
\newcommand{\docsubtitle}{WARDEN - Wireless Attack Reproduction and Detection Environment}

\usepackage{warden-style}

\begin{document}

% -- Portada -------------------------------------------------------------------
\wardencoverpage{0.1}{Abril 2026}{Borrador}

% -- Historial de Revisiones ---------------------------------------------------
\begin{wardenrevisions}
\revrow{0.1}{Abril 2026}{Santiago Valencia León}{Borrador inicial}
\revrow{0.1}{Abril 2026}{Daniel Yadir Perea Murillo}{Borrador inicial}
\revrow{0.1}{Abril 2026}{Sofia Valdez Londono}{Borrador inicial}
\end{wardenrevisions}

% -- Tabla de Contenidos -------------------------------------------------------
\tableofcontents
\newpage

% ==============================================================================
\section{Introducción}
% ==============================================================================

\subsection{Propósito}
Este documento detalla el modelo de amenaza (threat model) del subsistema
ofensivo de WARDEN, así como la narrativa, justificación técnica y reglas de
ejecución de la cadena de ataque que se reproduce en el laboratorio. Sirve
como complemento del Documento de Visión, especificando con rigor el lado
ofensivo del proyecto.

Un threat model es un artefacto estándar en proyectos de ciberseguridad que
formaliza, antes de la ejecución de cualquier ataque, las preguntas
fundamentales: \emph{quién} ataca, \emph{a quién} ataca, \emph{por qué} ataca,
\emph{cómo} ataca, \emph{qué logra} si tiene éxito, y \emph{qué reglas}
delimitan el ejercicio. La existencia de este documento garantiza que la
ejecución del subsistema ofensivo es deliberada, contenida y trazable.

\subsection{Audiencia}
Este documento está dirigido a:

\begin{itemize}
  \item Los integrantes del equipo de desarrollo, como referencia
    autoritativa de qué se permite ejecutar y qué no.
  \item El profesor evaluador, como evidencia de que el ejercicio ofensivo
    está formalmente planificado y delimitado.
  \item Cualquier evaluador externo o futuro contribuidor que necesite
    entender la lógica del lado ofensivo del proyecto.
\end{itemize}

\subsection{Naturaleza Académica}
Toda actividad descrita en este documento se realiza con fines estrictamente
educativos en el contexto de la asignatura Teoría de Sistemas de la
Universidad Tecnológica de Pereira. Las pruebas se ejecutan exclusivamente
contra equipos propiedad del equipo de desarrollo. Los datos manejados
durante las demostraciones son ficticios. La cadena de ataque descrita aquí
no se ejecuta, bajo ninguna circunstancia, contra redes, dispositivos o
usuarios ajenos al equipo.

% ==============================================================================
\section{Marco Conceptual}
% ==============================================================================

\subsection{Threat Modeling}
El modelado de amenazas (threat modeling) es la práctica de identificar,
documentar y analizar las amenazas a un sistema antes de su ejecución. En el
contexto de WARDEN, el sistema modelado no es un activo a defender, sino el
ejercicio mismo de ataque controlado. El threat model responde:

\begin{enumerate}
  \item ¿Qué actor malicioso se está simulando?
  \item ¿Cuáles son las capacidades, limitaciones y motivaciones de ese
    actor?
  \item ¿Qué activos son objetivo del ataque?
  \item ¿Qué técnicas emplea el actor para alcanzar sus objetivos?
  \item ¿Qué evidencias deja cada técnica?
  \item ¿Qué resultado esperado se considera ``éxito'' del atacante?
\end{enumerate}

\subsection{Cyber Kill Chain}
La cadena de ataque (kill chain) es un modelo conceptual desarrollado
originalmente por Lockheed Martin que describe las fases secuenciales de
una operación ofensiva. En WARDEN la cadena se compone de tres fases que se
ejecutan en orden y en las que el éxito de cada fase habilita la siguiente:

\begin{enumerate}
  \item \textbf{Reconocimiento y Distracción.} El atacante prepara el
    entorno y reduce la capacidad de la víctima de discernir señales
    legítimas de señales maliciosas.
  \item \textbf{Aislamiento.} El atacante desconecta a la víctima de su
    canal legítimo de comunicación.
  \item \textbf{Explotación.} El atacante captura las credenciales o el
    objetivo final, aprovechando el estado vulnerable de la víctima.
\end{enumerate}

\subsection{MITRE ATT\&CK}
MITRE ATT\&CK es una matriz pública mantenida por MITRE Corporation que
documenta tácticas y técnicas adversarias observadas en operaciones reales.
Cada fase de la cadena de WARDEN se mapea a técnicas específicas de esta
matriz. Este mapeo no es decorativo: ancla el ejercicio académico a un marco
profesional reconocido y permite que cualquier evaluador familiarizado con
la matriz identifique inmediatamente lo que el ejercicio reproduce.

% ==============================================================================
\section{Modelo de Atacante}
% ==============================================================================

\subsection{Perfil del Actor}
El atacante simulado por WARDEN corresponde al arquetipo del
\textbf{actor oportunista cercano} (proximity-based opportunistic attacker),
ampliamente documentado en la literatura de seguridad inalámbrica.

\begin{datatable}{L{3.5cm} L{12.3cm}}{Atributo & Caracterización}
\drow{Tipo de actor    & Atacante oportunista cercano. No tiene un objetivo individual específico predefinido; ataca a quien se encuentre conectado a la red que decide suplantar.}
\drow{Ubicación        & Físicamente cercano al objetivo, dentro del rango de cobertura WiFi. Típicamente opera desde espacios públicos con alta densidad de usuarios: cafeterías, salas de espera de aeropuerto, bibliotecas, campus universitarios.}
\drow{Capacidades      & Bajas a moderadas. Hardware accesible (microcontrolador económico). Conocimiento técnico de nivel intermedio sobre el protocolo 802.11. No requiere acceso previo a la infraestructura objetivo.}
\drow{Recursos         & Bajos. El equipamiento completo cuesta menos de 30 USD. No requiere infraestructura de comando y control.}
\drow{Persistencia     & Baja. Operación oportunista de corta duración (minutos a horas). No busca mantener acceso prolongado.}
\drow{Detectabilidad   & Moderada. La ejecución produce evidencias claras observables por sensores pasivos, pero típicamente no son monitoreados en redes domésticas o empresariales sin WIDS especializado.}
\drow{Motivación       & Robo de credenciales para reutilización en otros sistemas (correo institucional, redes sociales, banca, etc.).}
\end{datatable}

\subsection{Casos Reales de Referencia}
Este perfil de atacante y la cadena de técnicas descritas en este documento
no son una construcción teórica del equipo. Corresponden a operaciones
reales documentadas por la industria de seguridad:

\begin{itemize}
  \item \textbf{Campañas DarkHotel} (Kaspersky Lab, documentadas desde 2007).
    Operaciones contra ejecutivos hospedados en hoteles, donde los atacantes
    desplegaban Evil Twins sobre redes WiFi de hoteles para robar
    credenciales corporativas.
  \item \textbf{APT28 / Fancy Bear contra hoteles europeos} (FireEye, 2017).
    Uso documentado de deauthentication seguido de Evil Twin contra
    huéspedes específicos.
  \item \textbf{Demostraciones en DEF CON Wall of Sheep} (anuales). Pruebas
    de concepto que muestran ataques de Evil Twin contra asistentes a la
    conferencia.
  \item \textbf{Herramientas públicas como Wifiphisher} que automatizan
    exactamente la cadena de tres fases que reproduce WARDEN, evidencia de
    que esta combinación es el patrón de ataque establecido.
\end{itemize}

% ==============================================================================
\section{Activos Objetivo}
% ==============================================================================

\subsection{Inventario de Activos}
La cadena de ataque de WARDEN tiene los siguientes activos como objetivo,
todos propiedad del equipo de desarrollo:

\begin{datatable}{L{3.5cm} L{4.5cm} L{7.8cm}}{Activo & Identificador & Rol en el ataque}
\drow{Router de laboratorio & SSID: \texttt{LAB\_WARDEN\_UTP} \newline BSSID: por configurar & Red legítima a suplantar. Objeto del Evil Twin.}
\drow{Redmi Note 7 & MIUI 12.5 \newline MAC: por documentar & Víctima primaria. Objetivo del robo de credenciales.}
\drow{Dell Latitude E6220 & Ubuntu Server 22.04.5 LTS \newline MAC: por documentar & Víctima secundaria. Recibe Beacon Flood y Deauth.}
\end{datatable}

\subsection{Activos Ficticios para Demostración}
Durante la ejecución del ataque, el portal cautivo solicita ``credenciales''
que son ficticias y solo existen para fines demostrativos. Bajo ninguna
circunstancia se introducen credenciales reales en el portal cautivo de
WARDEN.

\textbf{Credenciales ficticias estándar para demostración:}

\begin{datatable}{L{3.5cm} L{5cm} L{7.3cm}}{Campo & Valor de demostración & Justificación}
\drow{Usuario  & \texttt{demo.user@warden.test} & Dominio inexistente, sufijo \texttt{.test} reservado por RFC 2606}
\drow{Contraseña & \texttt{demo-password-12345} & Cadena obviamente sintética, sin correspondencia con cuenta real}
\end{datatable}

% ==============================================================================
\section{Reglas de Ejecución (Rules of Engagement)}
% ==============================================================================

Las siguientes reglas son inviolables y delimitan el alcance permitido del
subsistema ofensivo:

\begin{enumerate}
  \item \textbf{Aislamiento de objetivos.} Los ataques se dirigen
    exclusivamente al BSSID del router de laboratorio y a las direcciones
    MAC de los dispositivos víctima del equipo. La ESP32 nunca opera con
    objetivos genéricos ni con direcciones broadcast a redes ajenas.
  \item \textbf{Solo datos ficticios.} El portal cautivo solo acepta y
    registra credenciales ficticias. Si por error un miembro del equipo
    ingresa una credencial real, se elimina inmediatamente del log.
  \item \textbf{Sin internet.} La infraestructura de laboratorio nunca está
    conectada a internet durante las pruebas. Esto previene cualquier
    fuga accidental de tráfico capturado.
  \item \textbf{Ubicación controlada.} Las pruebas se realizan en espacios
    con baja densidad de redes WiFi vecinas. Antes de cada sesión, se
    verifica el espectro inalámbrico circundante.
  \item \textbf{Duración limitada.} Cada sesión de prueba tiene una duración
    máxima planificada. La cadena de ataque no se ejecuta de forma continua
    o desatendida.
  \item \textbf{Registro obligatorio.} Cada ejecución se documenta en una
    bitácora de sesión con: fecha, hora de inicio y fin, participantes
    presentes, fases ejecutadas, resultados observados.
  \item \textbf{Detención inmediata.} Cualquier evidencia de que un
    dispositivo o red ajena está siendo afectada por la ejecución implica
    la detención inmediata del subsistema ofensivo.
  \item \textbf{No despliegue fuera del laboratorio.} El firmware del
    atacante nunca se ejecuta fuera de las sesiones de laboratorio
    autorizadas. La ESP32 con el firmware ofensivo se mantiene desconectada
    cuando no está en uso.
  \item \textbf{Distribución del firmware.} El firmware del subsistema
    ofensivo se publica únicamente con propósitos académicos y se
    acompaña siempre de su documentación, incluyendo este threat model.
  \item \textbf{Confirmación ética obligatoria en el Panel del Atacante.}
    El Panel del Atacante exige al Operador marcar explícitamente una
    casilla de confirmación ética antes de habilitar el botón de
    lanzamiento del ataque. Esta confirmación humana se complementa con
    la validación automática del Validador Ético en el firmware, que
    se ejecuta tanto al cambiar la configuración como al recibir la
    petición \texttt{POST /attack/start}. Una configuración rechazada
    retorna HTTP 403 sin iniciar ningún ataque.
  \item \textbf{Aislamiento de la red de control.} La red WiFi
    \texttt{WARDEN\_CONTROL} servida por la ESP32 para que la laptop
    del Operador acceda a la API REST está protegida con WPA2 y
    credenciales conocidas únicamente por el equipo. La red de control
    no tiene salida a internet y opera exclusivamente como canal
    administrativo entre el Operador y la ESP32.
\end{enumerate}

% ==============================================================================
\section{Cadena de Ataque}
% ==============================================================================

La cadena WARDEN se compone de tres fases secuenciales que reproducen el
patrón \emph{Reconocimiento → Aislamiento → Explotación}.

\subsection{Diagrama de la Cadena}

\vspace{0.3cm}
\begin{verbatim}
+----------------------------------------------------------+
|  FASE 1: BEACON FLOODING                                 |
|  Objetivo: saturar el espectro y reducir confianza       |
|  MITRE ATT&CK: T1498 - Network Denial of Service         |
|  Duracion tipica: 30-60 segundos                         |
+----------------------------------------------------------+
                          |
                          v
+----------------------------------------------------------+
|  FASE 2: DEAUTHENTICATION                                |
|  Objetivo: desconectar a la victima del AP legitimo      |
|  MITRE ATT&CK: T1499 - Endpoint Denial of Service        |
|  Duracion tipica: 5-15 segundos                          |
+----------------------------------------------------------+
                          |
                          v
+----------------------------------------------------------+
|  FASE 3: EVIL TWIN + PORTAL CAUTIVO                      |
|  Objetivo: capturar credenciales por suplantacion        |
|  MITRE ATT&CK: T1557 - Adversary-in-the-Middle           |
|                T1056.003 - Web Portal Capture            |
|  Duracion tipica: hasta que la victima ingresa datos     |
+----------------------------------------------------------+
                          |
                          v
                  CADENA COMPLETADA
\end{verbatim}

\subsection{Fase 1: Beacon Flooding}

\begin{atktable}
\atkrow{Fase}{1 - Reconocimiento y Distracción}
\atkrow{Nombre del ataque}{Beacon Flooding}
\atkrow{Técnica MITRE}{T1498 - Network Denial of Service}
\atkrow{Objetivo táctico}{Saturar el entorno inalámbrico de la víctima con redes falsas, reduciendo la capacidad del usuario y de las herramientas automatizadas para discernir entre redes legítimas y maliciosas.}
\atkrow{Mecanismo}{La ESP32 emite continuamente beacons 802.11 con SSIDs aleatorios y BSSIDs falsificados. Cada beacon anuncia una ``red disponible'' inexistente.}
\atkrow{Volumen}{Al menos 50 beacons únicos por segundo durante la duración de la fase.}
\atkrow{Direccionalidad}{Broadcast en el canal del laboratorio. No se dirige a un cliente específico.}
\atkrow{Efecto sobre la víctima}{La lista de redes WiFi disponibles del dispositivo se llena de SSIDs sospechosos, dificultando identificar la red real entre el ruido.}
\atkrow{Evidencia generada}{Múltiples beacons con BSSIDs no vistos previamente, generados a tasa anómala desde un solo emisor físico (la ESP32).}
\atkrow{Detección esperada}{El subsistema defensivo detecta una tasa anómala de beacons únicos por segundo (D-01).}
\atkrow{Caso real referenciado}{Demostraciones en DEF CON con \texttt{mdk3} / \texttt{mdk4} y herramientas equivalentes.}
\end{atktable}

\subsection{Fase 2: Deauthentication}

\begin{atktable}
\atkrow{Fase}{2 - Aislamiento}
\atkrow{Nombre del ataque}{Deauthentication}
\atkrow{Técnica MITRE}{T1499 - Endpoint Denial of Service}
\atkrow{Objetivo táctico}{Desconectar a la víctima de la red legítima para forzarla a reconectarse, momento en el que su dispositivo es vulnerable a conectarse al Evil Twin que se desplegará en la fase 3.}
\atkrow{Mecanismo}{La ESP32 envía frames de deautenticación 802.11 con dirección de origen falsificada (impersonando al AP legítimo) dirigidos a la MAC del cliente víctima, o broadcast al BSSID legítimo.}
\atkrow{Volumen}{Suficiente para mantener desconectada a la víctima durante el periodo en que el Evil Twin está activo.}
\atkrow{Direccionalidad}{Dirigido al BSSID del router legítimo y a las MAC de los dispositivos víctima del equipo. Nunca a otros BSSIDs.}
\atkrow{Efecto sobre la víctima}{El cliente recibe el deauth, asume que el AP lo está desconectando, y procede a buscar redes para reconectarse. En Android (Redmi Note 7), esto sucede en menos de 5 segundos. En Linux (Dell), el evento queda registrado en \texttt{dmesg} y \texttt{journalctl}.}
\atkrow{Evidencia generada}{Frames de deauth con código de razón típico (1, 4, 7 según el caso), originados desde una MAC que suplanta al AP legítimo.}
\atkrow{Detección esperada}{El subsistema defensivo detecta deauths dirigidos al BSSID protegido (D-02).}
\atkrow{Caso real referenciado}{APT28 contra hoteles europeos (FireEye, 2017). Herramienta \texttt{aireplay-ng} en pentests.}
\end{atktable}

\subsection{Fase 3: Evil Twin con Portal Cautivo Falso}

\begin{atktable}
\atkrow{Fase}{3 - Explotación}
\atkrow{Nombre del ataque}{Evil Twin con Portal Cautivo Falso}
\atkrow{Técnica MITRE}{T1557 - Adversary-in-the-Middle \newline T1056.003 - Web Portal Capture}
\atkrow{Objetivo táctico}{Hacer que la víctima se conecte a un AP controlado por el atacante (Evil Twin) y capturar credenciales mediante una página web falsa servida como portal cautivo.}
\atkrow{Mecanismo}{La ESP32 cambia de modo emisor de beacons falsos a modo Access Point completo. Anuncia un SSID idéntico al de la red legítima (\texttt{LAB\_WARDEN\_UTP}) pero como red abierta (sin contraseña). Sirve un servidor DHCP, un servidor DNS que redirige todas las consultas a la propia ESP32, y un servidor web que presenta el portal cautivo falso.}
\atkrow{Reconexión de la víctima}{El Redmi Note 7, recién desconectado por la Fase 2, busca el SSID conocido. Al encontrarlo abierto, intenta conectarse. Android detecta automáticamente que es una red con portal cautivo y lo presenta al usuario.}
\atkrow{Portal cautivo}{Página web mínima diseñada por el equipo. Solicita ``re-autenticación'' con campos de usuario y contraseña. Identificada visiblemente como componente de demostración educativa de WARDEN.}
\atkrow{Direccionalidad}{El AP solo es alcanzable físicamente por dispositivos en el rango de la ESP32. Los servicios sirven exclusivamente a clientes que se conectan voluntariamente.}
\atkrow{Efecto sobre la víctima}{La víctima ve el portal cautivo. En el laboratorio, el operador de la víctima ingresa credenciales ficticias estándar (\texttt{demo.user@warden.test} / \texttt{demo-password-12345}).}
\atkrow{Captura}{Las credenciales ficticias se registran por consola serial de la ESP32 y, simultáneamente, se publican en el endpoint \texttt{GET /credentials} de la API REST para que el Operador del Atacante las visualice en el Panel del Atacante. El almacenamiento es exclusivamente en memoria volátil.}
\atkrow{Evidencia generada}{Aparición de un nuevo BSSID anunciando un SSID idéntico a uno conocido, pero con configuración de seguridad distinta (abierto vs. WPA2).}
\atkrow{Detección esperada}{El subsistema defensivo detecta el SSID duplicado con BSSID distinto al protegido (D-03).}
\atkrow{Caso real referenciado}{Campañas DarkHotel (Kaspersky). Herramienta \texttt{wifiphisher}.}
\end{atktable}

% ==============================================================================
\section{Resumen de Detección Esperada}
% ==============================================================================

Para cada fase de la cadena ofensiva, el subsistema defensivo de WARDEN
debe identificar evidencias específicas. La siguiente tabla resume la
correspondencia entre el lado ofensivo y el defensivo:

\begin{datatable}{L{2cm} L{4.5cm} L{4.5cm} L{4.3cm}}{Fase & Ataque & Firma de Detección & Objetivo de Detector}
\drow{1 & Beacon Flooding & Tasa anómala de beacons únicos por segundo & D-01}
\drow{2 & Deauthentication & Frames de deauth dirigidos al BSSID protegido & D-02}
\drow{3 & Evil Twin & SSID duplicado con BSSID distinto al BSSID protegido & D-03}
\drow{Cadena & Las tres fases secuenciales & Correlación temporal de las tres firmas en intervalo corto & D-07}
\end{datatable}

% ==============================================================================
\section{Variantes y Escenarios de Demostración}
% ==============================================================================

\subsection{Escenario A: Cadena Completa Automática}
La ESP32 ejecuta las tres fases en secuencia automática al presionar un
botón o recibir un comando por consola serial. Esta modalidad demuestra el
ataque tal como ocurriría en el mundo real.

\subsection{Escenario B: Ejecución Fase por Fase Manual}
El operador del atacante selecciona qué fase ejecutar y cuándo. Permite
estudiar cada fase en aislamiento y mostrar al detector señales más
limpias, útil para presentaciones didácticas.

\subsection{Escenario C: Solo Detección sobre PCAP de Referencia}
El subsistema defensivo se ejecuta sobre archivos PCAP previamente
capturados durante sesiones autorizadas. Útil cuando el subsistema ofensivo
no está disponible en una demostración.

\end{document}
